\scanpage{130}
\mspage{4}
% archive sheet number: 129
On peut \unsure{donc|aussi} dire que la suite \deleted{transfinie}\inserted{ordinale} des factorisations fonctorielles $\unsure{\Phi|\mathcal F}_\alpha(f)$ de $f$ variable, se déduit de la factorisation \inserted{fonctorielle} $\Phi_1(f)$ de $f$
\diagram{\begin{tikzcd}[ampersand replacement=\&]
X \arrow[r,"i_1(f)"] \arrow[rr,bend left=20,"f"] \& X_1(f) \arrow[r,"p_1(f)"] \& Y
\end{tikzcd}}
par ``itération transfinie à droite'', i.e. en regardant les opérations $\Phi_\alpha$ comme des foncteurs
\[
p_\alpha\deleted{\unsure{\varphi|f}}\left\{\begin{aligned}
\underline{\mathrm{Fl}}(M)&\longrightarrow\underline{\mathrm{Fl}}(M)\\
f&\longmapsto p_\alpha(f)
\end{aligned}\right.
\]
\margin{\uline{NB} $p_\alpha$ respecte les structures de Catégorie sur $M$, via le foncteur $\operatorname{but}:\underline{\mathrm{Fl}}(M)\longrightarrow M$. i.e. $\underline{\mathrm{Fl}}(M)$ est considéré comme objet de $\mathrm{CAT}_{/M}$}
\unsure{munies|munis} de \deleted{\unsure{un foncteur}} hom. fonctoriels (au dessus de $M$)
\[\mathrm{id}_{\mathrm{Fl}\,M}\longrightarrow p_\alpha\deleted{\unsure{\Phi|\varphi}}.\]

Ainsi, si on considère
\[
\left\{\begin{aligned}
&\underline{\mathrm{Fl}}\,M\text{ comme objet de }\mathrm{CAT}_{/M}\\
&\qquad\text{par le foncteur }f\longmapsto\operatorname{but}f\\
&\text{où }\mathrm{CAT}\text{ est la catégorie des Catégories}\\
&\qquad\text{\uline{éléments de} }\unsure{\mathcal V|V|\mathcal U},\text{ univers tel que}\\
&M\in\mathcal V\text{ d'où }\left.\begin{aligned}\mathrm{Cat}^{(\unsure{\mathrm{ens}|\mathrm{et}})}&=\mathrm{Cat}^{\unsure{\mathcal U|U}}\\M,\mathrm{Fl}\,M\deleted{^{\unsure{\mathcal U|*}}}&\end{aligned}\right\}\in\mathrm{CAT}
\end{aligned}\right.
\]
on a un endofoncteur $p_1$ de $\underline{\mathrm{Fl}}(M)$ sur $M$, et un hom. fonctoriel
\[i_1:\mathrm{id}_{\underline{\mathrm{Fl}}(M)_{/M}}\longrightarrow p_1\qquad\text{i.e.}\tag*{1.13.}\]
\diagram{\begin{tikzcd}[ampersand replacement=\&]
X \arrow[rr,"f"] \arrow[dr,"i_1(f)"'] \& \& Y \\
\& \begin{matrix}X_1(f)\\\Vert\\\operatorname{source}p_1(f)\end{matrix} \arrow[ur,"p_1(f)"'] \&
\end{tikzcd}}
, et la description de la suite ordinale des endofoncteurs $p_\alpha$ de $\underline{\mathrm{Fl}}\,M_{/M}$ résulte de cette seule donnée, compte tenu que les petites limites ind. filtrées existent dans les fibres de $\underline{\mathrm{Fl}}\,M\xrightarrow{\operatorname{but}}M$, \deleted{\unsure{et}}
